\scanpage{36}
que cet élément a une image \uline{non nulle} dans $Z_p$, donc que l'extension ne provient pas d'une extension de groupes $p$-divisibles \uline{sur $A$}.

\uline{Question} Quel est l'élément obtenu de $Z_p$ ! est-ce tjrs un entier, resp. un entier $>0$ ?

2º) \uline{Cas} $\operatorname{car}K\unsure{>|\ne}0$ \uline{donc} $\operatorname{car}K=p$

\deleted{\unsure{Cela|Ce cas}} Il y a lieu d'étudier $T_p(K)$, qui ne s'explicite plus en général exclusivement en termes de modules galoisiens.

A) Cas $j$ entier (bonne réduction). a) \unsure{Si|si} $j \bmod \unsure{\mathfrak m}$ est ordinaire, alors $E_K$ est ordinaire, donc
\figure{\begin{tikzpicture}\fill (0,0) circle (0.025);\draw[line width=0.6pt] (0.07,0)--(1.2,0);\fill (0.27,0) ellipse (0.08 and 0.025);\fill (1.2,0) ellipse (0.05 and 0.025);\end{tikzpicture}}
$T_p(E_K)$ est une extension d'un $L_K$ par $L_K^{\vee}\otimes T_p(G_m)$, et l'extension est donnée par un invariant dans $L^{-2}\otimes A^{**}$ (pas \unsure{nulle|nul} dans $L^{-2}\otimes$
\margin{(\unsure{si|qui} $k$ alg. clos)}
\figure{\begin{tikzpicture}\draw[->] (0,-0.35) .. controls (-0.1,0.1) and (0.2,0.15) .. (1.5,0.3) -- (4.3,0.65);\end{tikzpicture}}
\margin{35}
